\documentclass[a4paper,10pt]{article}

%%%%%%%%%%%% PREÁMBULO %%%%%%%%%%%%%%%%%%%%%

% Paquetes

\usepackage[utf8]{inputenc}
\usepackage[spanish]{babel}
\usepackage[T1]{fontenc}
\usepackage{array}
\usepackage{amsmath}
\usepackage{geometry}
\geometry{margin=2.5cm}
\usepackage{listings}
\usepackage{graphicx}
\usepackage{xcolor}

% Comandos

% Opciones

% Estilo
\lstdefinelanguage{Circom}{
	keywords={template, signal, input, output, component},
	keywordstyle=\color{blue}\bfseries,
	morecomment=[l]{//},   % comentarios //
	morecomment=[s]{/*}{*/}, % comentarios /* */
	morestring=[b]",       % cadenas con comillas
}

\lstset{
	basicstyle=\ttfamily\small,
	keywordstyle=\color{blue},
	commentstyle=\color{green!50!black},
	stringstyle=\color{red!60!black},
	numbers=left,
	numberstyle=\tiny,
	stepnumber=1,
	numbersep=10pt,
	showstringspaces=false,
	tabsize=2,
	breaklines=true
}

%%%%%%%%%%%%%%%%%%%%%%%%%%%%%%%%%%%%%%%%%%%%%%

\begin{document}
	%Introducción al documento
	\title{Síntesis, traducción y verificación de sistemas de ecuaciones utilizados en Protocolos de Conocimiento}
	\author{Clara Rodríguez, Alejandro Hernández y Lucía Martínez}
	\date{\today}  
	\maketitle 
	
	\begin{abstract}
		\begin{center}
			Este documento está destinado para la transcripción de operadores C a su equivalente en lenguaje circom.
			
		\end{center}
	\end{abstract}
	
	\section*{Guía de contenidos}
	
\end{document}